\section{Discussion}\label{sec:discussion}
%
Iteration on unbounded containers is an
\emph{out-of-gas vulnerability} of smart contracts~\cite{AlbertGRS19}.
We provided a sound static boundedness analysis
based on abstract interpretation.
Boundedness analysis might also be tackled via
cost analysis~\cite{AlbertAGP11} of the contracts~\cite{PerezKLMH20},
after translation
into Datalog (MadMax~\cite{GrechKJBSS18}, in general unsound)
or Horn clauses and polyhedra (GASTAP/GASOL~\cite{AlbertCGRR20}, sound):
if a loop has constant complexity, then it must be bounded.
We have seen that this seems to fail in practice
(Sec.~\ref{sec:implementation}) since the two problems are distinct.

%We strived to make this paper self-contained. Therefore,
Our boundedness
analysis builds solely on a simple, relatively weak side-effect analysis
%(Sec.~\ref{sec:side_effect})
that distinguishes lists by their elements' type.
This can be limiting: the
analysis cannot distinguish two lists with elements of the same type and
would mark both as unbounded when either is unbounded.
This can lead to imprecision, that can be overcome by identifying
lists in a more precise way, for instance, by a creation point analysis,
a concretization of class analysis~\cite{TipP00}.

We expect our analysis to be sound, as typical with abstract
interpretation, but we still miss a concrete semantics and a proof of
correctness for the abstract operations (Fig.~\ref{fig:bo_derivations_commands}),
that remain as future work.
